\documentclass[10pt,a4paper]{amsart}
\usepackage[margin=19mm]{geometry}
\usepackage{amsmath,amssymb,mathtools}
\usepackage[scr=rsfs,cal=euler]{mathalfa}
\usepackage{xfrac}
\usepackage[unicode,hidelinks,bookmarks=false]{hyperref}
\newcommand{\ie}{i.e.\ }
\newcommand{\cf}{cf.\ }
\newcommand{\resp}{respectively\ }
\newcommand{\Subsection}{\subsection}
\providecommand{\nobreakdash}{\nobreak\hbox{-}}
\setlength{\emergencystretch}{2em}
\multlinegap0pt
\usepackage{bm}
\usepackage{bbm}
\usepackage{dsfont}
\usepackage{xspace}
\usepackage{cancel}
\usepackage{mathtools}
\usepackage{amsthm}
\makeatletter
\@ifundefined{thm}{\newtheorem{thm}{Theorem}}{}
\@ifundefined{lem}{\newtheorem{lem}[thm]{Lemma}}{}
\makeatother
\newcommand{\FF}{\mathbb{F}}
\newcommand{\Ical}{\mathbf{I}}
\newcommand{\NN}{\mathbb{N}}
\newcommand{\dep}{\operatorname{dep}}
\newcommand{\ww}[1]{\mathfrak{#1}} 
\renewcommand{\H}{\mathbf{H}}
\title[On $u$-Multiple Zeta Values in Positive Characteristic]{On $u$-Multiple Zeta Values in Positive Characteristic}
\author{Hung-Chun Tsui}
\date{Source: arXiv:2604.03618v1 (CC BY 4.0)}
\begin{document}
\maketitle

\noindent\emph{Source and licence.} On $u$-Multiple Zeta Values in Positive Characteristic. Version arXiv:2604.03618v1 (CC BY 4.0), distributed under the Creative Commons Attribution 4.0 International licence, as recorded in the arXiv licence metadata for this exact version and shown on the abstract page https://arxiv.org/abs/2604.03618v1. Full rights notice: licenses/arxiv-2604.03618-CC-BY-4.0.txt.\par\smallskip

\noindent\textit{Editorial scope.} Counted unit: the Thm whose source label is \texttt{thm:abstract-r-shuffle} in the article (source0.tex lines 906--910, source label \texttt{thm:abstract-r-shuffle}), delivered together with the article's own complete original proof of it (source0.tex lines 911--964, one \texttt{proof} environment of 54 lines). No mathematics is written, repaired or re-derived: statement and proof are the article's own text line for line. Required context retained uncounted: eq:df\_Delta (lines 367--370); lem:H<d=sumHd (lines 860--863); lem:d=d product <d (lines 865--868); lem:shuffle base (lines 870--875). Every definition, display and label of those lines is retained unchanged. Omitted from this package and not redistributed: 1--366, 371--859, 864--864, 869--869, 876--905, 965--1797. Nothing inside a retained excerpt is omitted or altered: every retained body is delivered exactly as the article's lines are written, so no omission receipt of any kind accompanies this package. Citations: the retained excerpts carry no citation command at all, so no bibliography entry of the article is transcribed and no reference is redistributed. Reference adaptation: none was needed and none was made; every reference inside the delivered unit resolves inside it, so no unresolved reference or citation is produced. Printed slips kept literal: no printed word, symbol, hypothesis or number of the counted statement or of its proof was altered. Layout and class: the article's own class is \texttt{amsart} with packages including silence, amsfonts, amssymb, amsmath, latexsym, fontenc, fouriernc, longtable, booktabs, tabularx, lipsum, amsthm, array, xcolor; the wrapper is the lane's pinned \texttt{amsart} editorial wrapper at 10pt on A4, which restates the theorem-like environments the retained text needs and adds the article's own macro definitions that the retained text needs (\texttt{\textbackslash FF}, \texttt{\textbackslash H}, \texttt{\textbackslash Ical}, \texttt{\textbackslash NN}, \texttt{\textbackslash dep}, \texttt{\textbackslash ww}), with extra wrapper packages (bm, bbm, dsfont, xspace, cancel, mathtools). The delivered numbering is the unit's own. Assets: no retained span contains a figure environment, a graphics-inclusion command, a PDF-page inclusion, a table or a TikZ picture; no image file is copied into this package, the package declares no asset dependency and no image file is redistributed with it. Licence: version arXiv:2604.03618v1 of this article is distributed under the Creative Commons Attribution 4.0 International licence, as recorded in the arXiv licence metadata for this exact version and shown on the abstract page https://arxiv.org/abs/2604.03618v1. A full rights notice is delivered at licenses/arxiv-2604.03618-CC-BY-4.0.txt. Characters: every retained excerpt is pure ASCII, so the delivered engine, XeLaTeX with its default Latin Modern text font, reports no missing character.
\begin{equation}\label{eq:df_Delta}
\Delta_{r_1,s_1}^{i,j}=\begin{cases} (-1)^{r_1-1}\binom{j-1}{r_1-1}+(-1)^{s_1-1}\binom{j-1}{s_1-1}, & \text{if } (r-1)\mid j \text{ and }j>0,\\ 0, & \text{otherwise},
    \end{cases}
\end{equation} 

\begin{lem}\label{lem:H<d=sumHd}
    Let $0\leq d<d_0$ and $\ww{s}=(s_1,\ldots,s_m)\in\Ical$ be a non-empty index. Then we have
    $$\H_{<d}(\ww{s})=\sum_{0\leq d_1<d}\H_{d_1}(\ww{s}).$$
\end{lem}

\begin{lem}\label{lem:d=d product <d}
     Let $0\leq d<d_0$ and $\ww{s}=(s_1,\ldots,s_m)\in\Ical$ be a non-empty index.  Then we have
    $$\H_{d}(\ww{s})=\H_{d}(s_1)\H_{<d}(\ww{s}^-).$$
\end{lem}

\begin{lem}\label{lem:shuffle base}
     Let $0\leq d< d_0$ and $r_1,s_1\in\NN$. Then we have
    $$\H_{d}(r_1)\H_{d}(s_1)=\H_{d}(r_1+s_1)+\sum_{\substack{i+j=r_1+s_1}}\Delta_{r_1,s_1}^{i,j} \H_{d}(i,j)$$
    where
    $\Delta_{r_1,s_1}^{i,j}$ is defined in \eqref{eq:df_Delta}.
\end{lem}

\begin{thm}\label{thm:abstract-r-shuffle}
    Let $1\leq d\leq d_0$. Then $\widehat{\H_{<d}}$ is an $\FF_p$-algebra homomorphism. That is, 
    $$\widehat{\H_{<d}}(x_{\ww{r}}\ast x_{\ww{s}})=\widehat{\H_{<d}}(x_{\ww{r}})\widehat{\H_{<d}}(x_{\ww{s}})=\H_{<d}(\ww{r})\H_{<d}(\ww{s})$$
    for any indices $\ww{r},\ww{s}\in\Ical$.
\end{thm}

\begin{proof}
     We proceed by induction on the total depth $k=\dep(\ww{r})+\dep(\ww{s})$. When $\ww{r}$ or $\ww{s}$ is empty, the results is clear. In particular, the base case $k=1$ holds. Suppose the result holds for all $k<\ell$. We may assume both $\ww{r}$ and $\ww{s}$ are non-empty.  Let 
    $\ww{r}=(r_1,\ldots,r_n),\ww{s}=(s_1,\ldots,s_m)\in\Ical$ be non-empty indices
    with $m+n=\ell$.
    Then by Lemma \ref{lem:H<d=sumHd}, we have
    \begin{align*}
        &\widehat{\H_{<d}}(x_{\ww{r}})\widehat{\H_{<d}}(x_{\ww{s}})\\
        &=\left(\sum_{0\leq d_1<d}\H_{d_1}(\ww{r})\right)\left(\sum_{0\leq d_2<d}\H_{d_2}(\ww{s})\right)
        \\
        ={} &\begin{multlined}[t]
        \sum_{0\leq d_2<d_1<d}\H_{d_1}(\ww{r})\H_{d_2}(\ww{s})+\sum_{0\leq d_1<d_2<d}\H_{d_1}(\ww{r})\H_{d_2}(\ww{s})+\sum_{0\leq d_2=d_1<d}\H_{d_1}(\ww{r})\H_{d_2}(\ww{s})\\
        =\sum_{0\leq d_1<d}\H_{d_1}(\ww{r})\left(\sum_{0\leq d_2<d_1}\H_{d_2}(\ww{s})\right)+\sum_{0\leq d_2<d}\left(\sum_{0\leq d_1<d_2}\H_{d_1}(\ww{r})\right)\H_{d_2}(\ww{s})  
            \\+ \sum_{0\leq d_3<d}\H_{d_3}(\ww{r})\H_{d_3}(\ww{s})
        \end{multlined}\\
        &=\sum_{0\leq d_1<d}\H_{d_1}(\ww{r})\H_{<d_1}(\ww{s})+\sum_{0\leq d_2<d}\H_{<d_2}(\ww{r})\H_{d_2}(\ww{s})  
            + \sum_{0\leq d_3<d}\H_{d_3}(\ww{r})\H_{d_3}(\ww{s})
    \end{align*}
    For the third equality, note that empty sums are taken to be zero.
    Hence, by Lemma \ref{lem:d=d product <d}, we obtain
    \begin{multline}\label{eq:shuffle2}
\widehat{\H_{<d}}(x_{\ww{r}})\widehat{\H_{<d}}(x_{\ww{s}})
= \sum_{0\leq d_1<d} \H_{d_1}(r_1) \H_{<d_1}(\ww{r}^-) \H_{<d_1}(\ww{s}) \\
+ \sum_{0\leq d_2<d} \H_{d_2}(s_1) \H_{<d_2}(\ww{s}^-) \H_{<d_2}(\ww{r}) \\
+ \sum_{0\leq d_3<d} \H_{d_3}(s_1) \H_{d_3}(r_1) \H_{<d_3}(\ww{r}^-) \H_{<d_3}(\ww{s}^-).
\end{multline}
By induction hypothesis and Lemma \ref{lem:d=d product <d}, the first summand in \eqref{eq:shuffle2} becomes
    \begin{align*}
        \sum_{0\leq d_1<d}\H_{d_1}({r_1})\H_{<d_1}({\ww{r}^-})\H_{<d_1}(\ww{s})&=\sum_{0\leq d_1<d}\H_{d_1}(r_1)\widehat{\H_{<d_1}}(x_{\ww{r}^-}\ast x_{\ww{s}})\\
        &=
        \sum_{0\leq d_1<d}\widehat{\H_{d_1}}(x_{r_1}(x_{\ww{r}^-}\ast x_{\ww{s}}))
        \\
        &=\widehat{\H_{<d}}(x_{r_1}(x_{\ww{r}^-}\ast x_{\ww{s}})).
    \end{align*}
The second equality can be checked immediately by writing $x_{\ww{r}^-}\ast x_\ww{s}=\sum_{i}c_ix_{\ww{t}_i}$, $c_i\in\FF_p$ and using the $\FF_p$-linearity. Similarly, the second summand in \eqref{eq:shuffle2} equals to
    $$\sum_{0\leq d_2<d} \H_{d_2}(s_1) \H_{<d_2}(\ww{s}^-) \H_{<d_2}(\ww{r})=\widehat{\H_{<d}}(x_{s_1}(x_{\ww{r}}\ast x_{\ww{s}^-})).$$
Finally, by Lemmas \ref{lem:H<d=sumHd}, \ref{lem:d=d product <d}, \ref{lem:shuffle base} and induction hypothesis, the third summand in \eqref{eq:shuffle2} is
\begin{align*}
    &\sum_{0\leq d_3<d} \H_{d_3}(s_1) \H_{d_3}(r_1) \H_{<d_3}(\ww{r}^-) \H_{<d_3}(\ww{s}^-)\\
    &=\sum_{0\leq d_3<d}\H_{d_3}(s_1) \H_{d_3}(r_1) \widehat{\H_{<d_3}}(x_{\ww{r}^-}\ast x_{\ww{s}^-}) \\
    &=
    \sum_{0\leq d_3<d}\left(\H_{d_3}(r_1+s_1)+\sum_{{i+j=r_1+s_1}}\Delta_{r_1,s_1}^{i,j} \H_{d_3}(i,j)\right)\widehat{\H_{<d_3}}(x_{\ww{r}^-}\ast x_{\ww{s}^-})\\
    &=
    \sum_{0\leq d_3<d}\H_{d_3}(r_1+s_1)\widehat{\H_{<d_3}}(x_{\ww{r}^-}\ast x_{\ww{s}^-}) +\sum_{0\leq d_3<d}\sum_{{i+j=r_1+s_1}}\Delta_{r_1,s_1}^{i,j} \H_{d_3}(i,j)\widehat{\H_{<d_3}}(x_{\ww{r}^-}\ast x_{\ww{s}^-}) \\
    &=
    \widehat{\H_{<d}}(x_{r_1+s_1}(x_{\ww{r}^-}\ast x_{\ww{s}^-})) +\sum_{{i+j=r_1+s_1}}\Delta_{r_1,s_1}^{i,j}\sum_{0\leq d_3<d}\H_{d_3}(i)\widehat{\H_{<d_3}}(x_{\ww{r}^-}\ast x_{\ww{s}^-}) \H_{<d_3}(j)\\
    &=
    \widehat{\H_{<d}}(x_{r_1+s_1}(x_{\ww{r}^-}\ast x_{\ww{s}^-})) +\sum_{{i+j=r_1+s_1}}\Delta_{r_1,s_1}^{i,j}\sum_{0\leq d_3<d}\H_{d_3}(i)\widehat{\H_{<d_3}}(x_{\ww{r}^-}\ast x_{\ww{s}^-}) \widehat{\H_{<d_3}}(x_j)\\
    &=
    \widehat{\H_{<d}}(x_{r_1+s_1}(x_{\ww{r}^-}\ast x_{\ww{s}^-})) +\sum_{{i+j=r_1+s_1}}\Delta_{r_1,s_1}^{i,j}\sum_{0\leq d_3<d}\H_{d_3}(i)\widehat{\H_{<d_3}}((x_{\ww{r}^-}\ast x_{\ww{s}^-})\ast x_j)\\
    &=
    \widehat{\H_{<d}}(x_{r_1+s_1}(x_{\ww{r}^-}\ast x_{\ww{s}^-})) +\sum_{{i+j=r_1+s_1}}\Delta_{r_1,s_1}^{i,j}\widehat{\H_{<d}}(x_i((x_{\ww{r}^-}\ast x_{\ww{s}^-})\ast x_j)).
\end{align*}
This completes the proof.
\end{proof}

\end{document}
